\chapter{Chapter 6 --- Conclusions and Planned
Work}\label{chapter-6-conclusions-and-planned-work}

This report covers the first two phases of the project. This chapter
states what they establish, what they do not, and how the remaining
phases are designed --- including two commitments made here, in advance,
rather than after the data are seen.

\section{6.1 What Phases 1 and 2
establish}\label{what-phases-1-and-2-establish}

\subsection{6.1.1 The two research questions
answered}\label{the-two-research-questions-answered}

\textbf{RQ1 --- How much do SZZ variants disagree, and in which
direction?} Substantially, and structurally. No variant exceeds
\textbf{27.2\% precision} against the reference. Variants agree with one
another at \ensuremath{\kappa} up to \textbf{0.933} while agreeing with the reference at \ensuremath{\kappa}
\textbf{0.158--0.202} --- four to six times more strongly with each
other than with the thing they are estimating. Bias is bifurcated: B-SZZ
over-flags (\ensuremath{\rho}\ensuremath{_0} = 0.263), while L-SZZ and R-SZZ under-flag (\ensuremath{\rho}\ensuremath{_1} = 0.733
and 0.700). \textbf{Two decades of refinement have relocated error
rather than reduced it}, moving it from the false-positive column to the
false-negative column.

\textbf{RQ2 --- How does the label source shift measured performance
across regimes?} It shifts it significantly --- 0.034 to 0.064 MCC under
realistic streaming, against a deployable baseline near 0.097.
\textbf{But the evaluation protocol shifts it far more.} Random
\emph{k}-fold inflates a high-capacity model by \textbf{+0.137 MCC};
scoring against the training labels inflates it by a further
\textbf{+0.230}. The two protocol effects together are roughly nine
times the label-source effect, and both scale with model capacity.

\subsection{6.1.2 The finding that reframes the
project}\label{the-finding-that-reframes-the-project}

The project set out to measure what SZZ noise costs. It found that
\textbf{the cost of SZZ noise is small relative to the cost of the
field's evaluation conventions} --- and that the two are not
independent.

Self-scoring is only possible \emph{because} the labels are heuristic.
If the field had ground truth, it would score against it, and the +0.230
gap could not exist. \textbf{The labelling failure creates the
evaluation failure.} A result reported under random \emph{k}-fold with
self-scored SZZ labels is not a measurement of defect-prediction
capability, and on this corpus the difference between that protocol and
a defensible one is 0.41 against 0.10 MCC.

\subsection{6.1.3 Practical significance, stated
honestly}\label{practical-significance-stated-honestly}

The measurement critique is statistically overwhelming --- rank-biserial
+1.000 on the leakage effect, 21 of 21 projects, surviving global
correction. \textbf{Practical significance is more modest and should be
reported as such.} Even with the reference labels and an honest
protocol, deployable MCC on this corpus is approximately \textbf{0.097}.

This report does not claim that better labels would make JIT-SDP good.
It claims the field's reported numbers overstate what it currently
achieves by roughly a factor of four. The low ceiling is itself a
finding: a field reporting 0.41 and delivering 0.10 has a calibration
problem independent of whether 0.10 is useful.

\subsection{6.1.4 The reconciliation a reader will otherwise trip
over}\label{the-reconciliation-a-reader-will-otherwise-trip-over}

Studies that evaluate SZZ against developer-informed oracles report
precision in the range 0.39--0.66. This report reports 0.186--0.272.
\textbf{The two are not in conflict, and the difference is the
denominator.}

Those studies score \emph{conditionally} on bug-fixing commits already
known to have an inducing commit: given a defect, which earlier commit
introduced it? This report scores over an entire project history of
27,319 commits, the overwhelming majority of which are clean: is
\emph{this} commit, drawn from all of history, defect-introducing? A
variant right two times in three on the first question can be wrong four
times in five on the second. \textbf{Any comparison between the two
literatures has to reconcile the denominator first.}

\section{6.2 What Phases 1 and 2 do not
establish}\label{what-phases-1-and-2-do-not-establish}

Three questions are raised by this report and are not answerable from
its evidence. Each is stated with the design that would answer it.

\subsection{6.2.1 Which SZZ error actually costs the
learner}\label{which-szz-error-actually-costs-the-learner}

Phase 2 shows that label quality matters under streaming. It does
\textbf{not} show whether the cost comes from the false positives SZZ
invents or from the defects it misses. The simplest explanation is ruled
out: the penalty does not track a variant's recall (\ensuremath{\rho} = \ensuremath{-}0.18), and
L-SZZ and B-SZZ have almost identical gaps despite recalls of 0.267 and
0.641.

The question cannot be settled by comparing label sources, because whole
labelling procedures differ in many ways at once. \textbf{It needs an
intervention.} Phase 3's design takes real B-SZZ labels and uses the
reference to correct \emph{one error type at a time} --- removing the
false positives, or restoring the false negatives --- holding the
stream, the learner and the latency model fixed. The difference between
conditions is then attributable to the repair.

\textbf{Two controls must be built into that design from the start},
because this report has already identified the confounds they would
otherwise create:

\begin{itemize}
\tightlist
\item
  \textbf{Error matching.} B-SZZ carries 6,565 false positives against
  837 false negatives, an 8:1 imbalance (\S{}4.2.1). ``Repair all of each''
  therefore conflates \emph{which error is costlier per label} with
  \emph{which error there is more of}. The design must include a
  matched-mass condition.
\item
  \textbf{Delivery.} A restored false negative has no linked fix and so
  no natural arrival time. If arrival times are imputed from the
  empirical latency distribution, a null result cannot distinguish
  ``these labels do not matter'' from ``these labels cannot arrive in
  time to matter''. The design must include an immediate-delivery
  ceiling condition.
\end{itemize}

\subsection{6.2.2 Why an online learner should be sensitive at
all}\label{why-an-online-learner-should-be-sensitive-at-all}

ORB handles class imbalance by oversampling at a rate set from the
\emph{observed} rate of defect labels --- that is, from label
\textbf{quantity}. It has no access to label \textbf{quality}, so a
false positive receives the same amplification as a true one.

That makes amplification a plausible route by which SZZ noise damages an
online learner, and it makes a measurable prediction: the oversampling
rate applied to defect-labelled arrivals should be an inverse function
of how many such labels the stream delivers, \emph{regardless of whether
they are correct}. \textbf{This report does not test that.}
Instrumenting the learner's internal rate across a calibrated dose grid
would.

\subsection{6.2.3 Whether the latency effect can be
isolated}\label{whether-the-latency-effect-can-be-isolated}

\S{}5.3 withdraws a decomposition of the batch-to-streaming drop because
the contrast was not identified, and reports the replacement 2\ensuremath{\times}2 as
unresolved --- every contrast, including the interaction, has an
interval spanning zero at \emph{n} = 21.

The obstacle is that switching label sources changes the learner's
training signal and its evaluation window together. A design that
\textbf{holds the learner fixed by construction} --- injecting
controlled noise into a single label source rather than switching
between sources --- removes that confound. It also requires a true
no-latency control: an arm in which \emph{every} label, positive and
negative, arrives immediately. A fixed-delay arm is not a no-latency
arm, since it delays the 91.5\% majority class as heavily as the
minority.

\section{6.3 Commitments made in
advance}\label{commitments-made-in-advance}

Two are recorded here so that they precede the data rather than follow
it.

\textbf{The mitigation phase will be pre-registered.} Hypotheses,
primary model, acceptance bar and multiplicity correction will be
committed to version control before the first run, and the commit
history will timestamp the order. Phases 1 and 2 are exploratory and are
labelled as such throughout; the global Holm column is what makes their
headline claims defensible without a registration.

\textbf{The non-degradation gate is the acceptance bar, not the headline
test.} A mitigation that improves performance on noisy labels while
damaging it on clean ones is not adoptable, because a practitioner
cannot know in advance which case they are in. That gate will be applied
before any effect size is examined.

\section{6.4 Threats to validity}\label{threats-to-validity}

\textbf{The reference is not ground truth.} It is \texttt{git\ blame}
seeded with human-verified fix lines (\S{}3.2.2). Comparisons against it
isolate the cost of tangled commits, not total labelling error, and
\textbf{all measured noise is therefore a lower bound}: blame error is
present on both sides of every comparison and cancels.

\textbf{The reference condition's label timing is SZZ-derived.} This is
the most serious threat in the study. The published labels carry no
native fix-to-inducing linkage, so arrival times were reconstructed from
the union of SZZ mappings. The reference is blame-derived in
construction \emph{and} SZZ-derived in timing. A coverage sensitivity
analysis (67.8\% \ensuremath{\rightarrow} 100\%) bounds part of the exposure; the timing itself
is not bounded, and every streaming result inherits it.

\textbf{The corpus is SZZ-shaped.} Changes adding no new lines were
excluded by the dataset authors, so defects introduced purely by
deletion cannot appear. Because the reference is itself blame output,
blame-unreachability is not an available explanation for any measured
false negative.

\textbf{Project heterogeneity is measured and does not drive the
results.} Projects span 544--4,026 commits and 19--335 defects and are
weighted equally. Excluding those below a 25-defect floor moves no
headline estimate by more than 0.002 MCC, and all headline claims still
survive correction (\S{}5.4).

\textbf{Effects below roughly 0.04 MCC are not resolvable at \emph{n} =
21.} This is load-bearing for the latency decomposition, which is
reported as unresolved rather than as null.

\textbf{Nothing here is pre-registered}, which is why global
multiplicity correction is reported alongside the within-family
correction, and why the headline claims are those surviving the global
column.

\textbf{One benchmark family.} Twenty-one Apache Java repositories of a
single dataset lineage.

\section{6.5 Closing}\label{closing}

The report began by asking how much SZZ label noise costs a defect
prediction model. The answer --- between a third and two-thirds of the
available signal under realistic streaming --- turned out to be less
interesting than a question it exposed along the way: \textbf{how much
of what the field reports is measurement artefact rather than
capability?}

On this corpus, most of it. A number that reads as 0.41 under the
field's standard protocol is 0.10 once the model is scored against
something other than the heuristic that trained it and evaluated in an
order that time permits. That gap --- roughly four-fifths of the
reported signal --- is quantified here, bounded by confidence intervals,
corrected for multiple comparisons, stable across every analytical
choice tested, and reproducible from a single commit hash.

What the report cannot yet say is \emph{why}: which of SZZ's two errors
an online learner cannot survive, and whether a learner can be built to
withstand it. Those are the next two phases, and \S{}6.2 states both the
designs and the controls that this report has already shown they will
need.
